\subsection{注记与进一步结果}
\label{sec:ch1-notes-further-results}

\subsubsection{参考资料}
\label{subsec:ch1-references}

三角级数方面的经典参考书是 Zygmund 的著作 \MBUnresolvedReference{citation}{[21]}, 它也是随后几章中各项结果的重要参考书. 不过, 这部著作有时不易查阅. Bary 的著作 \MBUnresolvedReference{citation}{[1]} 是关于三角级数的另一部综合性参考书.

\hypertarget{chapter-01-page-032}{}
Katznelson \MBUnresolvedReference{citation}{[10]} 以及 Dym 与 McKean \MBUnresolvedReference{citation}{[4]} 对 Fourier 级数与积分有出色的论述. R. E. Edwards 的著作 \MBUnresolvedReference{citation}{[5]} 从较现代的观点对 Fourier 级数进行了详尽研究. 还推荐 Weiss 的论文 \MBUnresolvedReference{citation}{[20]} 和 Körner 的著作 \MBUnresolvedReference{citation}{[12]}. J. P. Kahane 关于 Fourier 级数及其对数学概念发展之影响的优秀历史论述, 见 \MBUnresolvedReference{citation}{[9]} 的前半部. E. González-Velasco 的著作 \emph{Fourier Analysis and Boundary Value Problems}(Academic Press, New York, 1995)包含许多将 Fourier 分离变量法应用于偏微分方程的例子, 也包含历史资料. 另见同一作者的 \emph{Connections in mathematical analysis: the case of Fourier series}, Amer. Math. Monthly 99 (1992), 427--441. O. G. Jørsboe 与 L. Melbro 的著作 \emph{The Carleson--Hunt Theorem on Fourier Series}, Lecture Notes in Math. 911, Springer-Verlag, Berlin, 1982, 专门讨论该定理的证明. 其原始文献是 L. Carleson 的论文 \emph{On convergence and growth of partial sums of Fourier series}, Acta Math. 116 (1966), 135--157, 以及 R. Hunt 的论文 \emph{On the convergence of Fourier series}, Orthogonal Expansions and their Continuous Analogues (Proc. Conf., Edwardsville, Ill., 1967), pp. 235--255, Southern Illinois Univ. Press, Carbondale, 1968. Kolmogorov 关于一个 Fourier 级数处处发散的 $L^1$ 函数的例子, 出现在 \emph{Une série de Fourier-Lebesgue divergente partout}(C. R. Acad. Sci. Paris 183 (1926), 1327--1328)中.

\subsubsection{多重 Fourier 级数}
\label{subsec:ch1-multiple-fourier-series}

令 $\mathbb T^n$ 为 $n$ 维环面(可将它与商群 $\mathbb R^n/\mathbb Z^n$ 等同). 定义在 $\mathbb T^n$ 上的函数等价于定义在 $\mathbb R^n$ 上并且关于每个变量都以 $1$ 为周期的函数. 若 $f\in L^1(\mathbb T^n)$, 则可定义其 Fourier 系数为
\[
 \widehat f(\nu)=\int_{\mathbb T^n}f(x)e^{-2\pi ix\cdot\nu}\,dx,
 \qquad \nu\in\mathbb Z^n,
\]
并用这些系数构造 $f$ 的 Fourier 级数
\[
 \sum_{\nu\in\mathbb Z^n}\widehat f(\nu)e^{2\pi ix\cdot\nu}.
\]
可以证明若干与单变量 Fourier 级数类似的结果, 但随着 $n$ 增大, 所需的正则性条件越来越严格. 见 Stein 与 Weiss \MBUnresolvedReference{citation}{[18, Chapter 7]}.

\subsubsection{Poisson 求和公式}
\label{subsec:ch1-poisson-summation-formula}

设函数 $f$ 满足: 对某个 $\delta>0$,
\[
 |f(x)|\le A(1+|x|)^{-n-\delta},
 \qquad
 |\widehat f(\xi)|\le A(1+|\xi|)^{-n-\delta}.
\]
特别地, $f$ 和 $\widehat f$ 都连续. 则
\hypertarget{chapter-01-page-033}{}
\[
 \sum_{\nu\in\mathbb Z^n}f(x+\nu)
 =\sum_{\nu\in\mathbb Z^n}\widehat f(\nu)e^{2\pi ix\cdot\nu}.
\]
这个等式(更准确地说, $x=0$ 的情形)称为 Poisson 求和公式, 它不过是反演公式. 左端定义了 $\mathbb T^n$ 上的函数, 其 Fourier 系数恰为 $\widehat f(\nu)$.

\subsubsection{Gibbs 现象}
\label{subsec:ch1-gibbs-phenomenon}

在 $(-1/2,1/2)$ 上令 $f(x)=\operatorname{sgn}(x)$. 例如由 Dirichlet 判别法可知, 对所有 $x$, $S_Nf(x)$ 收敛到 $f(x)$. 在 $0$ 的右侧, 部分和在 $1$ 附近振荡; 但与可能的预期相反, 当 $N$ 增大时, 它们超过 $1$ 的量并不趋于 $0$. 可以证明
\[
 \lim_{N\to\infty}\sup_xS_Nf(x)
 =\frac2\pi\int_0^\pi\frac{\sin y}{y}\,dy
 \approx1.17898\ldots.
\]
只要函数存在跳跃间断点, 就会出现这种现象. 它以 J. Gibbs 命名; Gibbs 于 1899 年在 Nature 59 上宣布了这一现象, 但 H. Wilbraham 早在 1848 年已经发现它. 见 Dym 与 McKean \MBUnresolvedReference{citation}{[4, Chapter 1]}, 以及 E. Hewitt 与 R. E. Hewitt 的论文 \emph{The Gibbs--Wilbraham phenomenon: an episode in Fourier analysis}, Arch. Hist. Exact Sci. 21 (1979/80), 129--160.

以 Cesàro 可和性代替逐点收敛可以消除 Gibbs 现象. 若 $m\le f(x)\le M$, 则由 \eqref{eq:ch1-fejer-kernel-properties} 中 Fejér 核的前两个性质, $m\le\sigma_Nf(x)\le M$. 事实上, 可以证明, 若在区间 $(a,b)$ 上 $m\le f(x)\le M$, 则对任意 $\varepsilon>0$, 当 $N$ 充分大时, 在 $(a+\varepsilon,b-\varepsilon)$ 上有
\[
 m-\varepsilon\le\sigma_Nf(x)\le M+\varepsilon.
\]

\subsubsection{Hausdorff--Young 不等式}
\label{subsec:ch1-sharp-hausdorff-young}

推论~\ref{cor:ch1-hausdorff-young} 是 Riesz--Thorin 插值的直接应用. 但实际上有更强的不等式: 若 $1\le p\le2$, 则
\[
 \|\widehat f\|_{p'}
 \le\left(\frac{p^{1/p}}{(p')^{1/p'}}\right)^{n/2}\|f\|_p.
\]
这个不等式是锐的, 因为对 $f(x)=e^{-\pi|x|^2}$ 等号成立. 这一结果由 W. Beckner 证明(\emph{Inequalities in Fourier analysis}, Ann. of Math. 102 (1975), 159--182); $p$ 为偶数的特殊情形更早由 K. I. Babenko 证明(\emph{An inequality in the theory of Fourier integrals (Russian)}, Izv. Akad. Nauk SSSR Ser. Mat. 25 (1961), 531--542).

\hypertarget{chapter-01-page-034}{}
在同一篇论文中, Beckner 还证明了 Young 不等式(推论~\ref{cor:ch1-young-convolution})的锐化形式.

\subsubsection{$L^2(\mathbb R)$ 中 Fourier 变换的特征函数}
\label{subsec:ch1-fourier-eigenfunctions}

由于 Fourier 变换的周期为 $4$, 若 $f$ 满足 $\widehat f=\lambda f$, 则必有 $\lambda^4=1$. 因此, $\lambda=\pm1,\pm i$ 是 Fourier 变换仅有的可能特征值. 引理~\ref{lem:ch1-gaussian-fourier-transform} 表明 $\exp(-\pi x^2)$ 是属于特征值 $1$ 的特征函数. Hermite 函数给出其余特征函数: 对 $n\ge0$,
\[
 h_n(x)=\frac{(-1)^n}{n!}\exp(\pi x^2)
 \frac{d^n}{dx^n}\exp(-\pi x^2)
\]
满足 $\widehat h_n=(-i)^nh_n$. 将这些函数归一化,
\[
 e_n=\frac{h_n}{\|h_n\|_2}
 =\left[(4\pi)^{-n}\sqrt{2n!}\right]^{1/2}h_n,
\]
便得到 $L^2(\mathbb R)$ 的一组正交归一基, 并且
\[
 \widehat f=\sum_{n=0}^{\infty}(-i)^n\langle f,e_n\rangle e_n.
\]
于是 $L^2(\mathbb R)$ 分解为直和 $H_0\oplus H_1\oplus H_2\oplus H_3$, 其中 Fourier 变换在子空间 $H_j$, $0\le j\le3$, 上的作用是将函数乘以 $i^j$.

这种定义 $L^2(\mathbb R)$ 上 Fourier 变换的方法归功于 N. Wiener, 可见其著作 \emph{The Fourier Integral and Certain of its Applications}, original edition, 1933; Cambridge Univ. Press, Cambridge, 1988. 另见 Dym 与 McKean \MBUnresolvedReference{citation}{[4, Chapter 2]}.

在高维情形, Fourier 变换的特征函数是 Hermite 函数的乘积, 每个坐标变量对应一个因子. 另见 \MBUnresolvedReference{section}{第 4 章第 7.2 节}.

\subsubsection{解析算子族的插值}
\label{subsec:ch1-analytic-family-interpolation}

Riesz--Thorin 插值定理有一个由 E. M. Stein 给出的有力推广(见 Stein 与 Weiss \MBUnresolvedReference{citation}{[18, Chapter 5]}). 令
\[
 S=\{z\in\mathbb C:0\le\operatorname{Re}z\le1\},
\]
并令 $\{T_z\}_{z\in S}$ 为一族算子. 若给定两个函数 $f,g\in L^1(\mathbb R^n)$, 映射
\[
 z\longmapsto\int_{\mathbb R^n}T_z(f)g\,dx
\]
在 $S$ 的内部解析, 在边界上连续, 并且存在常数 $a<\pi$, 使
\[
 e^{-a|\operatorname{Im}z|}
 \log\left|\int_{\mathbb R^n}T_z(f)g\,dx\right|
\]
对所有 $z\in S$ 一致有界, 则称这个算子族是可容许的.

\hypertarget{chapter-01-page-035}{}
\begin{Theorem}
\label{thm:ch1-stein-interpolation}
设 $\{T_z\}$ 为可容许算子族, 并假设对 $1\le p_0,p_1,q_0,q_1\le\infty$ 和 $y\in\mathbb R$,
\[
 \|T_{iy}f\|_{q_0}\le M_0(y)\|f\|_{p_0},
 \qquad
 \|T_{1+iy}f\|_{q_1}\le M_1(y)\|f\|_{p_1},
\]
其中对某个 $b<\pi$,
\[
 \sup_{y\in\mathbb R}e^{-b|y|}\log M_j(y)<\infty,
 \qquad j=1,2.
\]
则当 $0<\theta<1$, $\operatorname{Re}z=\theta$, 且 $p,q$ 如定理~\ref{thm:ch1-riesz-thorin} 中定义时, 存在常数 $M_\theta$, 使得
\[
 \|T_zf\|_q\le M_\theta\|f\|_p.
\]
\end{Theorem}

\subsubsection{有限测度的 Fourier 变换}
\label{subsec:ch1-finite-measure-fourier-transforms}

如上所述, 若 $\mu$ 是有限 Borel 测度, 则 $\widehat\mu$ 是有界连续函数. 由这种方式得到的所有函数组成的集合可由下述结果刻画.

\begin{Theorem}
\label{thm:ch1-bochner}
若 $h$ 是有界连续函数, 则下列条件等价:
\begin{enumerate}[label=(\arabic*)]
\item 对某个正的有限 Borel 测度 $\mu$, 有 $h=\widehat\mu$;
\item $h$ 是正定的: 对任意 $f\in L^1(\mathbb R^n)$,
\[
 \int_{\mathbb R^n}\int_{\mathbb R^n}
 h(x-y)f(x)\overline{f(y)}\,dx\,dy\ge0.
\]
\end{enumerate}
\end{Theorem}

这一定理归功于 S. Bochner(\emph{Lectures on Fourier Integrals}, Princeton Univ. Press, Princeton, 1959; translated from \emph{Vorlesungen über Fouriersche Integrale}, Akad. Verlag, Leipzig, 1932). 另见 Katznelson \MBUnresolvedReference{citation}{[10, Chapter 6]}.

